% 14_body.tex — Day3 산출물 4/5 (⑭) 품질 통제 기록 — 결함 누적 · 감리 지적 대응 (빈 템플릿 / 완성 예시 공용 본문)
% 드라이버: 14_품질통제기록_빈템플릿.tex (\wbblanktrue) / 14_품질통제기록_완성예시.tex (\wbblankfalse)
% 내용 출처: PM트랙/Day3_워크북.md §4.3(항목 가이드) §4.4(빈 템플릿) §4.5(온담 P1 완성 예시본)
\documentclass[10pt,a4paper]{article}
\usepackage[a4paper,margin=16mm,top=14mm,bottom=20mm]{geometry}
\def\DocNum{4}
\def\DocTitle{품질 통제 기록 — 결함 누적 · 감리 지적 대응}
\def\DocSub{⑭ Quality Control Record (Defect Trend \& Audit Findings Response)}
\def\DocVariant{\B{빈 템플릿}{온담 P1 완성 예시}}
\usepackage{wbstyle}
\begin{document}

\docheader
 {\Bg{[정식 명칭과 코드]}{차세대 주문·재고 통합 플랫폼 구축 (ONE ONDAM P1)}}
 {\Bg{[기간 / 데이터 일자 / 발행]}{2026-02-02 ~ 09-30 (34주) / 발행 2026-10-05}}
 {\Bg{[PMO 품질 담당]}{P1 PMO 품질 담당 (이재현·정하늘 대조)\rd{7mm}}}
 {\Bg{[P1 PM]}{P1 PM (확인)\rd{7mm}}}

% ------------------------------------------------------------------ 1
\secbar{1. 문서 정보와 측정 정의}
\guide{결함의 정의(무엇을 세는가 — 요구 변경·해석 차이는 결함이 아니라 ⑪ 변경요청), 심각도 1~4의 정의(영향으로 — 우선순위와 섞지 않는다), 결함밀도 두 계층의 분자·분모(격자 인용), 도구·ID 체계, 감리 지적 등급 정의. 정의가 없으면 요구 변경이 결함으로 들어와 수행사 무상 작업 유인이 되고, 분모가 문서마다 달라진다.}
{\footnotesize
\begin{tabularx}{\linewidth}{|L{34mm}|Y|}\hline
\hd{항목} & \hd{내용} \\ \hline
\B{\rd{8mm}}{}기간 / 데이터 일자 & \Bg{[시작(근거) ~ 데이터 일자 / 발행, 부속 관계]}{2026-02-02(S3 설계 검토 시작) ~ 2026-09-30 / 발행 10-05, 성과보고 제9호 부속} \\ \hline
\B{\rd{12mm}}{}결함 정의 & \Bg{[무엇을 세는가 / 제외 / 판정 다툼의 결정자]}{시험(설계 검토회·단위·스프린트 통합·발주사 검증)에서 발견된 “기대 결과와의 불일치”로 결함 관리 도구에 등록·수행사가 인정한 건. \textbf{요구 변경·해석 차이는 결함이 아니라 변경요청(⑪)} — 판정 다툼은 P1 PMO 품질 담당이 3영업일 내 결정(수행사 이의 시 변경협의체). 중복·재현 불가는 제외} \\ \hline
\B{\rd{12mm}}{}심각도 & \Bg{[1 = … / 2 = … / 3 = … / 4 = … — 우선순위는 별도 열]}{(09-28 기준서 v1.0, 감리 5번 대응) \textbf{1} = 규제 요건 위반·데이터 손상·결제/재고 원장 중단 / \textbf{2} = 주요 기능 불가·우회 없음(Must 요구 미충족 포함) / \textbf{3} = 기능 오류이나 우회 있음 / \textbf{4} = 경미·UI·문구. 우선순위(P1~P3)는 별도 열} \\ \hline
\B{\rd{9mm}}{}밀도 — 작업패키지 계층 & \Bg{[분자 ÷ 분모, 기준 (격자), 초과 시 조치]}{스프린트 시험 결함(설계 검토 제외) ÷ 인수 FP 누적, 기준 \textbf{0.13건/FP}(격자, 수행사 PM 허용범위) — 초과 시 1단계 에스컬레이션} \\ \hline
\B{\rd{9mm}}{}밀도 — 프로젝트 계층 & \Bg{[분자 ÷ 분모, 기준, 판정 시점]}{통합시험(A14)~UAT(A16) 발견 결함 ÷ FP 기준선(BL-02 12,585), 기준 \textbf{0.4건/FP} + G3 심각도 1·2 미해결 0 — 판정 시점 2026-12-07 이후} \\ \hline
\B{\rd{8mm}}{}도구 · ID & \Bg{[결함 / 감리 지적 / 보안 취약점 번호 체계]}{결함 D-nnnn(도구 자동 채번), 감리 지적 AF-1-nn(1차), 보안 취약점 VA-1-nn(1회차, 결함 도구에 D-번호로 이중 등록)} \\ \hline
\B{\rd{10mm}}{}감리 등급 & \Bg{[중대 / 보통 / 경미 = 정의]}{\textbf{중대} = 게이트 통과 조건·규제·계약에 직접 영향, 시정 없이는 다음 단계 불가 / \textbf{보통} = 산출물 품질·추적성 결함, 기한 내 시정 / \textbf{경미} = 형식·문서 표기, 다음 감리 전 시정 (감리 계약 부속서 3)} \\ \hline
\B{\rd{8mm}}{}작성 / 검토 / 확인 & \Bg{[ ]}{P1 PMO 품질 담당 / 이재현(결함 데이터 대조), 정하늘(감리 부분) / P1 PM} \\ \hline
\end{tabularx}}

% ------------------------------------------------------------------ 2
\clearpage
\secbar{2. 주차별 결함 발견 · 해소 누적표}
\guide{주 단위 — 설계 검토 결함·시험 결함 발견, 누적 발견, 해소, 누적 해소, 미해결(backlog), 인수 FP 누적, 중간 밀도(인수 FP가 있는 주부터), \textbf{시험 노력(실행 TC)}. 월 단위로만 집계하면 평탄 구간이 안 보이고, 시험 노력이 없으면 평탄의 의미를 판독할 수 없다. 스프린트 경계를 표시한다.}
{\scriptsize\setlength{\tabcolsep}{2.5pt}
\begin{longtable}{|C{6mm}|L{11mm}|L{15mm}|C{11mm}|C{9mm}|C{11mm}|C{12mm}|C{11mm}|C{12mm}|C{11mm}|C{15mm}|C{13mm}|C{15mm}|}\hline
\hd{주} & \hd{시작일} & \hd{스프린트} & \hd{설계 검토} & \hd{시험} & \hd{주 발견} & \hd{누적 발견} & \hd{주 해소} & \hd{누적 해소} & \hd{미해결} & \hd{인수 FP 누적} & \hd{밀도 (시험÷FP)} & \hd{시험 노력 (TC 실행)} \\ \hline \endhead
\ifwbblank
\rd{3.6mm}1 & & & & & & & & & & & & \\ \hline
\rd{3.6mm}2 & & & & & & & & & & & & \\ \hline
\rd{3.6mm}3 & & & & & & & & & & & & \\ \hline
\rd{3.6mm}4 & & & & & & & & & & & & \\ \hline
\rd{3.6mm}5 & & & & & & & & & & & & \\ \hline
\rd{3.6mm}6 & & & & & & & & & & & & \\ \hline
\rd{3.6mm}7 & & & & & & & & & & & & \\ \hline
\rd{3.6mm}8 & & & & & & & & & & & & \\ \hline
\rd{3.6mm}9 & & & & & & & & & & & & \\ \hline
\rd{3.6mm}10 & & & & & & & & & & & & \\ \hline
\rd{3.6mm}11 & & & & & & & & & & & & \\ \hline
\rd{3.6mm}12 & & & & & & & & & & & & \\ \hline
\rd{3.6mm}13 & & & & & & & & & & & & \\ \hline
\rd{3.6mm}14 & & & & & & & & & & & & \\ \hline
\rd{3.6mm}15 & & & & & & & & & & & & \\ \hline
\rd{3.6mm}16 & & & & & & & & & & & & \\ \hline
\rd{3.6mm}17 & & & & & & & & & & & & \\ \hline
\rd{3.6mm}18 & & & & & & & & & & & & \\ \hline
\rd{3.6mm}19 & & & & & & & & & & & & \\ \hline
\rd{3.6mm}20 & & & & & & & & & & & & \\ \hline
\rd{3.6mm}21 & & & & & & & & & & & & \\ \hline
\rd{3.6mm}22 & & & & & & & & & & & & \\ \hline
\rd{3.6mm}23 & & & & & & & & & & & & \\ \hline
\rd{3.6mm}24 & & & & & & & & & & & & \\ \hline
\rd{3.6mm}25 & & & & & & & & & & & & \\ \hline
\rd{3.6mm}26 & & & & & & & & & & & & \\ \hline
\rd{3.6mm}27 & & & & & & & & & & & & \\ \hline
\rd{3.6mm}28 & & & & & & & & & & & & \\ \hline
\rd{3.6mm}29 & & & & & & & & & & & & \\ \hline
\rd{3.6mm}30 & & & & & & & & & & & & \\ \hline
\rd{3.6mm}31 & & & & & & & & & & & & \\ \hline
\rd{3.6mm}32 & & & & & & & & & & & & \\ \hline
\rd{3.6mm}33 & & & & & & & & & & & & \\ \hline
\rd{3.6mm}34 & & & & & & & & & & & & \\ \hline
\rd{3.6mm}\textbf{합} & & & & & & & & & & & & \\ \hline
\else
1 & 02-02 & S3 & 18 & — & 18 & 18 & 4 & 4 & 14 & 0 & — & 검토회 \\ \hline
2 & 02-09 & & 26 & — & 26 & 44 & 15 & 19 & 25 & 0 & — & 검토회 \\ \hline
3 & 02-16 & S4 & 31 & — & 31 & 75 & 22 & 41 & 34 & 0 & — & 검토회 \\ \hline
4 & 02-23 & & 34 & — & 34 & 109 & 30 & 71 & 38 & 0 & — & 검토회 \\ \hline
5 & 03-02 & S5 & 29 & — & 29 & 138 & 31 & 102 & 36 & 0 & — & 검토회 \\ \hline
6 & 03-09 & & 33 & — & 33 & 171 & 33 & 135 & 36 & 0 & — & 검토회 \\ \hline
7 & 03-16 & S6 & 24 & — & 24 & 195 & 30 & 165 & 30 & 0 & — & 검토회 \\ \hline
8 & 03-23 & & 15 & — & 15 & 210 & 26 & 191 & 19 & 0 & — & 검토회 \\ \hline
9 & 03-30 & S7 & — & 14 & 14 & 224 & 18 & 209 & 15 & 0 & — & 210 \\ \hline
10 & 04-06 & & — & 22 & 22 & 246 & 20 & 229 & 17 & 620 & 0.058 & 210 \\ \hline
11 & 04-13 & S8 & — & 31 & 31 & 277 & 26 & 255 & 22 & 620 & 0.108 & 210 \\ \hline
12 & 04-20 & & — & 38 & 38 & 315 & 33 & 288 & 27 & 1,240 & 0.085 & 210 \\ \hline
13 & 04-27 & S9 & — & 36 & 36 & 351 & 34 & 322 & 29 & 1,240 & 0.114 & 210 \\ \hline
14 & 05-04 & & — & 41 & 41 & 392 & 38 & 360 & 32 & 1,840 & 0.099 & 210 \\ \hline
15 & 05-11 & S10 (3주) & — & 40 & 40 & 432 & 37 & 397 & 35 & 1,840 & 0.121 & 210 \\ \hline
16 & 05-18 & & — & 47 & 47 & 479 & 42 & 439 & 40 & 1,840 & 0.146 & 210 \\ \hline
17 & 05-25 & & — & 33 & 33 & 512 & 34 & 473 & 39 & 2,320 & 0.130 & 210 \\ \hline
18 & 06-01 & S11 & — & 52 & 52 & 564 & 44 & 517 & 47 & 2,320 & 0.153 & 210 \\ \hline
19 & 06-08 & & — & 46 & 46 & 610 & 35 & 552 & 58 & 2,790 & 0.143 & 210 \\ \hline
20 & 06-15 & S12 & — & 50 & 50 & 660 & 43 & 595 & 65 & 2,790 & 0.161 & 210 \\ \hline
21 & 06-22 & & — & 48 & 48 & 708 & 40 & 635 & 73 & 3,350 & 0.149 & 210 \\ \hline
22 & 06-29 & S13 & — & 44 & 44 & 752 & 41 & 676 & 76 & 3,350 & 0.162 & 210 \\ \hline
23 & 07-06 & & — & 39 & 39 & 791 & 36 & 712 & 79 & 3,910 & 0.149 & 210 \\ \hline
24 & 07-13 & S14 & — & 46 & 46 & 837 & 39 & 751 & 86 & 3,910 & 0.160 & 210 \\ \hline
25 & 07-20 & & — & 44 & 44 & 881 & 37 & 788 & 93 & 4,510 & 0.149 & 210 \\ \hline
26 & 07-27 & S15 & — & 38 & 38 & 919 & 35 & 823 & 96 & 4,510 & 0.157 & 210 \\ \hline
27 & 08-03 & & — & \textbf{12} & 12 & 931 & 17 & 840 & 91 & 5,130 & 0.141 & \textbf{95} \\ \hline
28 & 08-10 & S16 & — & \textbf{9} & 9 & 940 & 10 & 850 & 90 & 5,130 & 0.142 & \textbf{95} \\ \hline
29 & 08-17 & & — & \textbf{11} & 11 & 951 & 11 & 861 & 90 & 5,750 & 0.129 & \textbf{95} \\ \hline
30 & 08-24 & S17 & — & 27 & 27 & 978 & 19 & 880 & 98 & 5,750 & 0.134 & 230 \\ \hline
31 & 08-31 & & — & 44 & 44 & 1,022 & 33 & 913 & 109 & 6,390 & 0.127 & 230 \\ \hline
32 & 09-07 & S18 (3주) & — & 52 & 52 & 1,074 & 41 & 954 & 120 & 6,390 & 0.135 & 230 \\ \hline
33 & 09-14 & & — & 38 & 38 & 1,112 & 39 & 993 & 119 & 6,390 & 0.141 & 230 \\ \hline
34 & 09-21 & & — & 28 & 28 & \textbf{1,140} & 31 & \textbf{1,024} & \textbf{116} & \textbf{7,020} & \textbf{0.132} & 230 \\ \hline
\textbf{합} & & & \textbf{210} & \textbf{930} & \textbf{1,140} & & \textbf{1,024} & & & & & \\ \hline
\fi
\end{longtable}}
\B{\small\g{월 합계(설계 검토 / 시험 / 해소)와 시험 노력 요약: \hrulefill}\par}{\small 시험 노력(주당 실행 TC, [추가 설정]): 4~7월 평균 210 / 8월 1~3주 평균 \textbf{95} / 8월 4주~9월 평균 230. 월 합계: 설계 검토 2월 109·3월 101 = 210; 시험 4월 127·5월 161·6월 240·7월 167·8월 103·9월 132 = 930. 해소 2월 71·3월 138·4월 113·5월 151·6월 203·7월 147·8월 90·9월 111 = 1,024. 검산 스크립트 day3\_defects.py.\par}

% ------------------------------------------------------------------ 3 (가로)
\landscapeon
\secbar{3. 결함 누적곡선 판독}
\guide{곡선의 구간별 형태(상승·평탄·재상승)와 원인, \textbf{평탄화의 두 의미 판별}(수렴형 vs 시험 정지형 — 시험 노력 대조), backlog 추이(발견 − 해소)와 해소 지연의 원인, 향후 곡선 전망. “평탄 = 수렴”으로 단정하지 않는다. 곡선을 그리고 판독을 쓰지 않으면 8월 평탄이 “안정화”로 보고된다.}
{\footnotesize
\begin{tabularx}{\linewidth}{|L{26mm}|L{30mm}|L{30mm}|L{22mm}|L{26mm}|Y|L{36mm}|}\hline
\hd{구간} & \hd{기간} & \hd{발견 추이 (주당)} & \hd{시험 노력} & \hd{평탄 유형} & \hd{원인 · 판독} & \hd{backlog} \\ \hline
\ifwbblank
\rd{14mm}① \g{[설계 검토]} & & & & \g{[수렴형 / 정지형 / —]} & & \g{[전 → 후, 해소 우세?]} \\ \hline
\rd{14mm}② \g{[스프린트 시험 상승]} & & & & & & \\ \hline
\rd{14mm}③ \g{[평탄]} & & & & & & \\ \hline
\rd{14mm}④ \g{[재상승]} & & & & & & \\ \hline
\else
① 설계 검토 & 02-02 ~ 03-29 (S3–S6) & 18 → 34 → 33 → 15 & 검토회 주 2회 일정 & \textbf{수렴형 평탄}(주 7~8) & 검토 대상 설계서 47종이 S5까지 대부분 검토 완료, 마지막 2주는 잔여·재검토 → 발견 감소가 노력 감소 없이 일어남 = 수렴 & 38 → 19 (해소 우세) \\ \hline
② 스프린트 시험 상승 & 03-30 ~ 07-31 (S7–S15) & 14 → 52 (6월 최고) → 38 & 주 210 TC & 상승 후 완만 & 인수 FP 증가와 비례. 6월 밀도 0.16 최고 — R-05 지연 후 S11·S12에 재작업 몰림 & 15 → \textbf{96} (해소 병목: 개발자가 R3 설계·규제 항목 병행) \\ \hline
③ 8월 평탄 & 08-03 ~ 08-23 (S16) & \textbf{12 → 9 → 11} & 주 \textbf{95} TC (−55\%) & \textbf{시험 정지형 평탄} & 시험 인력 6명 중 4명이 규제 검증(A07, 142개 항목 증적)에 이동. S16 인수 FP 620은 규제 항목(SRS-701~720 등 검증형)이라 FP당 결함이 적음. \textbf{곡선만 보면 품질 개선, 노력을 보면 시험 중단} & 96 → 90 (해소도 감소) \\ \hline
④ 9월 재상승 & 08-24 ~ 09-27 (S17–S18) & 27 → 44 → 52 → 38 → 28 & 주 230 TC & 재상승 후 하강 시작 & 8월 미시험분(S16 비규제 항목)이 S17에 몰림. 9/21 주 28은 S18 3주차 마무리 — 수렴인지 S19 착수 지연에 따른 시험 대상 부족인지 10월 2주 관측 필요 & 90 → \textbf{116} (최고) \\ \hline
\fi
\end{tabularx}}
\B{\small\g{전망 (도착률 근거, 프로젝트 전체 곡선의 두 번째 봉우리, 지금 봐야 할 곡선): \hrulefill}\par}{\small 전망: 스프린트 시험 결함의 도착률은 인수 FP에 비례하므로 R3(S19–S24) 기간에도 주당 35~50이 이어진다. 프로젝트 전체의 곡선은 통합시험(12월)에서 두 번째 봉우리(Rayleigh 형태의 후반)를 갖고 UAT에서 꼬리가 된다. \textbf{지금 봐야 할 것은 발견 곡선이 아니라 backlog 곡선}이다 — 6월부터 해소가 발견을 따라가지 못해 116건이 쌓였고, 그중 심각도 3·4가 109건이라 인수에는 지장이 없지만 R3 개발자의 시간을 잠식한다(A10 −10의 원인 중 하나).\par}

% ------------------------------------------------------------------ 4 (가로)
\secbar[60mm]{4. 결함밀도 중간 계산과 전망}
\guide{작업패키지 계층 밀도(시험 결함 ÷ 인수 FP)와 기준 대조, 월별 추이, 프로젝트 계층 밀도의 전망(범위로, 근거 붙여)과 분모 변경(BL-02) 명시. \textbf{0.132를 0.4와 비교해 “여유 있음”으로 적지 않는다}(계층 혼동). 밀도를 미해결 기준으로 계산하지 않는다.}
{\footnotesize
\begin{tabularx}{\linewidth}{|L{34mm}|Y|L{30mm}|L{32mm}|L{64mm}|}\hline
\hd{계층} & \hd{산식} & \hd{값} & \hd{기준} & \hd{판정} \\ \hline
\ifwbblank
\rd{9mm}작업패키지 (스프린트 시험) & \g{[시험 결함] ÷ [인수 FP]} & & \g{[0.13]} & \g{[내/초과 → 에스컬레이션]} \\ \hline
\rd{9mm}월별 추이 (누적) & \g{[월: 값 …]} & & & \g{[착시 여부]} \\ \hline
\rd{9mm}설계 검토 포함 (참고) & & & — & \\ \hline
\rd{9mm}프로젝트 계층 (전망) & \g{[스프린트 밀도] × [유출률] + [UAT 가정]} & \g{[범위]} & \g{[0.4 ÷ BL-0n FP]} & \\ \hline
\rd{9mm}절대 건수 전망 & \g{[밀도 × FP]} & & & \\ \hline
\else
작업패키지 (스프린트 시험) & 930 ÷ 7,020 FP & \textbf{0.132건/FP} & 0.13 (수행사 PM) & \textbf{초과 (+1.5\%)} → 1단계 에스컬레이션(이재현 → P1 PM, 09-28) \\ \hline
월별 추이 (누적) & 4월 0.085 / 5월 0.130 / 6월 0.162(최고, 6/29 주) / 7월 0.157 / 8월 0.129 / 9월 0.132 & & & 6~7월 초과 상태가 지속되었으나 \textbf{8월 시험 정지로 일시 회복처럼 보였음} — 8월 값 0.129는 착시 \\ \hline
설계 검토 포함 총 밀도 (참고) & 1,140 ÷ 7,020 & 0.162 & — & 참고 — 설계 결함은 FP 인수 전 발견이므로 분모 부정합 \\ \hline
프로젝트 계층 (전망) & 스프린트 밀도 0.132 × 통합시험 유출률 35\% [추가 설정, Kan DRE 가정] = 0.046 + UAT 발견(업무 시나리오·사용성) 약 0.04 [추가 설정] & \textbf{약 0.09건/FP} (범위 0.07~0.13) & 0.4 (통합시험~UAT ÷ 12,585) & 이내 전망. 단 R-19 실현(통합시험 압축) 시 유출률 50\% → 0.11, UAT 결함 몰림 → 심각도 2 미해결 0 조건이 병목 \\ \hline
절대 건수 전망 & 0.09 × 12,585 & 약 1,130건 (통합시험 580 + UAT 550) & 0.4 × 12,585 = 5,034 & 기준 대비 22\% — 기준 0.4는 넉넉하고 실질 제약은 “심각도 1·2 미해결 0” \\ \hline
\fi
\end{tabularx}}

% ------------------------------------------------------------------ 5 (가로)
\clearpage
\secbar{5. 심각도 1·2 미해결 현황과 나이}
\guide{미해결 결함의 심각도 분포, 심각도 2 이상 \textbf{전 건}의 ID·모듈(WBS)·발견일·나이(영업일)·담당·기한·우회 여부·연결. 나이를 재지 않으면 3개월 된 심각도 2 결함이 “미해결 7건”에 숨는다. 보안 취약점은 결함 도구에도 등록해 G3 조건에서 빠지지 않게 한다.}
\B{\small\g{미해결 [ ] = 심각도 1: [ ] / 2: [ ] / 3: [ ] / 4: [ ]. 나이: 전체 중앙값 [ ], 심각도 2 중앙값 [ ]·최장 [ ]}\par}{\small 미해결 116 = 심각도 1: \textbf{0} / 2: \textbf{7} / 3: 58 / 4: 51. 나이(발견일 → 데이터 일자, 영업일): 전체 중앙값 9, 심각도 2 중앙값 8·최장 16.\par}
{\footnotesize\setlength{\tabcolsep}{3pt}
\begin{tabularx}{\linewidth}{|L{14mm}|C{6mm}|L{30mm}|L{12mm}|C{9mm}|Y|L{30mm}|L{12mm}|L{22mm}|L{34mm}|}\hline
\hd{ID} & \hd{심} & \hd{모듈 (WBS)} & \hd{발견일} & \hd{나이} & \hd{내용} & \hd{담당} & \hd{기한} & \hd{우회} & \hd{연결} \\ \hline
\ifwbblank
\rd{8mm}\g{D-nnnn} & & & & \g{[영업일]} & & & & \g{[없음/있음]} & \g{[IS-·감리·R-·TC]} \\ \hline
\rd{8mm} & & & & & & & & & \\ \hline
\rd{8mm} & & & & & & & & & \\ \hline
\rd{8mm} & & & & & & & & & \\ \hline
\rd{8mm} & & & & & & & & & \\ \hline
\rd{8mm} & & & & & & & & & \\ \hline
\rd{8mm} & & & & & & & & & \\ \hline
\rd{8mm} & & & & & & & & & \\ \hline
\else
D-0966 & 2 & API GW 인증 (1.5.2) & 09-08 & 16 & 인증 토큰 만료 미검증 — 만료 토큰으로 재고 조회 API 호출 가능 (보안 High VA-1-01) & 수행 통합팀 + API GW 벤더 & 10-16 & 없음 & IS-08, 감리 6번, 격자 지속가능성 \\ \hline
D-0871 & 2 & 3PL 폴백 전환 (1.3.3) & 09-14 & 12 & 준실시간 → 폴백 자동 전환 시 재고 반영 지연 15분 초과(최대 23분) — TC-INV-037 실패 & 수행 인터페이스팀 & 10-09 & 없음 (수동 전환 시 정상) & R-16, StRS-032 \\ \hline
D-0902 & 2 & POS 반품 (1.1.2) & 09-16 & 10 & 반품 처리 시 멤버십 적립 취소 누락 → 적립 포인트 과다 & 수행 POS팀 & 10-09 & 없음 & StRS 멤버십 (HR-02 하위) \\ \hline
D-0915 & 2 & 개인정보 파기 (1.5.1) & 09-18 & 8 & 파기 주기 배치 실행 로그 미기록 — 증적 불가 (SRS-708) & 수행 통합팀 & 10-16 & 없음 & 감리 6번, TC-REG-098 \\ \hline
D-0931 & 2 & 엑셀 발주 업로드 (1.4.2) & 09-21 & 7 & 양식 5종 중 2종(구형 xls) 인코딩 오류로 수량 0 처리 (CR-03 확장분) & 수행 포털팀 & 10-09 & 없음 (양식 3종은 정상) & CR-03, UAT-FRN-031 \\ \hline
D-0944 & 2 & SAP 출고 확정 IF-SAP-12 (1.5.4) & 09-22 & 6 & 출고 확정 인터페이스 타임아웃(30초) 시 재시도 없이 실패 처리 → 발주 상태 불일치 & 수행 SAP팀 & 10-16 & 수동 재전송 & StRS-073, CR-09 파라미터 영향 \\ \hline
D-0958 & 2 & MDM 상품 중복 검출 (1.5.3) & 09-24 & 4 & 중복 검출 규칙이 브랜드 간 동일 명칭 상품을 오탐 → 정본 정제 결과 왜곡 & 수행 통합팀 & 10-23 & 수동 예외 처리 & R-10, 1.5.8 정본 판정 전 필수 \\ \hline
\fi
\end{tabularx}}
\B{\small\g{내부 기준(통합시험 착수 시 조건)과 G3 조건의 관계, 주간 고정 안건: \hrulefill}\par}{\small 내부 기준: 통합시험 착수(12-07) 시 심각도 2 미해결 0, 심각도 3 미해결 30 이하. G3 조건(심각도 1·2 미해결 0)은 UAT 종료 시점 판정 — 지금부터 “심각도 2의 나이 10영업일 초과 건”을 주간 PM 회의 고정 안건으로.\par}

% ------------------------------------------------------------------ 6 (가로)
\clearpage
\secbar{6. 감리 지적사항 대응 관리대장 \B{}{— 1차(중간), 감리 수행 09-07~09-18, 보고서 접수 09-25, 정하늘 수석감리원}}
\guide{지적 번호·감리 원문 요지·등급·감리 기재 귀속·\textbf{발주사 판정 귀속(발주사/수행사/공동)과 근거}(WBS 담당 열·RACI의 A·결정 기한 기록·의존 통지 기록)·조치·담당·기한·이의 여부·상태. 감리 보고서를 그대로 받아 “수행사 조치”로 넘기지 않는다 — 발주사 미결정이 원인인 지적을 수행사가 고칠 수 없고, 변경협의체에서 대기 비용의 근거가 된다. 이의는 사실이 아니라 \textbf{귀속}에 대해서만, 서면으로.}
{\scriptsize\setlength{\tabcolsep}{2pt}
\begin{longtable}{|L{11mm}|L{44mm}|C{8mm}|C{10mm}|L{16mm}|L{52mm}|L{42mm}|L{22mm}|L{14mm}|L{14mm}|L{14mm}|}\hline
\hd{\#} & \hd{지적 (감리 원문 요지)} & \hd{등급} & \hd{감리 기재 귀속} & \hd{발주사 판정 귀속} & \hd{근거} & \hd{조치} & \hd{담당} & \hd{기한} & \hd{이의} & \hd{상태} \\ \hline \endhead
\ifwbblank
\rd{12mm}AF-\g{n}-01 & & \g{[중대/보통/경미]} & \g{[발주/수행]} & \g{[발주/수행/공동]} & \g{[WBS 담당 열 · RACI A · 결정 기록 · 통지 기록]} & & & & \g{[서면 일자 / 없음]} & \\ \hline
\rd{12mm}-02 & & & & & & & & & & \\ \hline
\rd{12mm}-03 & & & & & & & & & & \\ \hline
\rd{12mm}-04 & & & & & & & & & & \\ \hline
\rd{12mm}-05 & & & & & & & & & & \\ \hline
\rd{12mm}-06 & & & & & & & & & & \\ \hline
\rd{12mm}-07 & & & & & & & & & & \\ \hline
\rd{12mm}-08 & & & & & & & & & & \\ \hline
\else
AF-1-01 & RTM 하위 분해율 88\%(StRS 26건 SyRS 미분해), 시험 연결 76\% — G2 기준 100\% 대비 미달 & 보통 & 수행사 & \textbf{공동} & 미분해 26건은 R3·R4 항목으로 스프린트 내 설계 대상(누락 아님) — 단 RTM 갱신 주기(동결 시)가 감리 시점과 어긋남은 발주사 PMO(1.6.3) 소관 & S22 설계 확정분 FZ-11 반영, RTM 갱신을 스프린트 리뷰마다로 변경(PMO) & P1 PMO(요구사항) + 수행 설계 리더 & 11-25 (FZ-11) & 없음 (성격 설명 회신) & 진행 \\ \hline
AF-1-02 & 인터페이스 설계문서 47본 중 9본 표준 양식(연동 주체·주기·항목·오류 처리) 미충족 — 오류 처리 절 누락 & 보통 & 수행사 & \textbf{수행사} & WBS 1.5.5 수행 담당, 표준 양식은 G1 확정(감리 합의) & 9본 보완 & 수행 인터페이스팀(신임 PL) & 10-30 & 없음 & 진행 \\ \hline
\textbf{AF-1-03} & \textbf{재고 원장 P2 설비 인터페이스 규격(IF-P2-01~04)에 설비측 확인 서명 부재 — “수행사 설계 미흡”} & \textbf{중대} & 수행사 & \textbf{발주사 (프로그램)} & 규격은 WBS 1.3.1 \textbf{발주 담당}(박세연·한동석 협의), P2 전달은 RACI 9행 \textbf{A = 프로그램 PMO장}. T0(04-24) 전 전달 완료(04-17 발송 기록). 설비 벤더(타케다)가 05-29 규격 변경 요청 2건을 보냈고 \textbf{프로그램 PMO의 회신이 09-30 현재 미결}(회신 요청 06-05·07-10·08-21 기록). 수행사는 미결 상태를 06-08·08-24 서면 통지 & 프로그램 PMO장 → P2 PM 규격 확정 회신 10-16, 확인 서명 첨부. 수행사는 회신 후 5영업일 내 규격서 v1.1 갱신 & \textbf{프로그램 PMO장} (P1 PM 촉구, 수행 인터페이스팀 후속) & 10-16 (회신) / 10-23 (갱신) & \textbf{이의 제기 09-29} (귀속) & 진행 — 회신 대기 \\ \hline
\textbf{AF-1-04} & \textbf{가맹 SW 배포 범위(1.1.6) 설계 미확정 — 비호환 32점 대응안이 배포 패키지 설계에 미반영, “수행사 설계 미흡”} & 보통 & 수행사 & \textbf{발주사 (영업본부)} & 1.1.6 조사·대응안 결정은 WBS \textbf{발주 담당}(오정근 A, RACI 11행). 대응안 3개(교체·임대·수기)는 04-30 조사 보고에 제시되었고 \textbf{결정은 부속서 협상 후로 유보(오정근 05-08 회신)} — 08-14 동의 71\% 후에도 미결. 수행사는 결정 전까지 패키지 설계 착수 불가(1.1.6 카드 외부 의존) & 오정근: 대응안 결정 G2(11-27)까지 — 미동의 매장 컷오버 후 2차 배포 옵션 포함. 수행사는 결정 후 패키지 설계 S25 & \textbf{오정근} (매장운영팀장, 수행 POS팀 후속) & 11-27 & \textbf{이의 제기 09-29} (귀속) & 진행 \\ \hline
AF-1-05 & 단위·스프린트 시험 결함의 심각도 분류 기준 문서 부재 — 발견자 주관 분류 & 경미 & 수행사 & \textbf{수행사} & 시험 계획(1.6.4 수행) 부속 누락 & 심각도 기준서 v1.0 작성·합의(09-28 완료), 기존 결함 116건 재분류(10-09) & 수행 시험팀 + P1 PMO 품질 & 10-09 & 없음 & 기준서 완료, 재분류 진행 \\ \hline
AF-1-06 & 개인정보 142개 항목 검증 증적 중 파기 주기 11건의 증적 형식이 REG-01 매핑표와 불일치(로그 미첨부) & 경미 & 수행사 & \textbf{수행사} & 검증 142/142 통과는 유효(최영주 확인), 증적 형식 문제. D-0915와 동일 원인 & 파기 배치 로그 기능 수정(D-0915) 후 증적 재첨부, REG-01 갱신 & 수행 통합팀 (최영주 확인) & 10-16 & 없음 & 진행 \\ \hline
AF-1-07 & 형상관리 — 발주사 저장소 동기화 주기 주 1회 → 일 1회 권고, 2023 내부감사 지적(소스 인수) 대응 관점 & 경미 & 발주사 & \textbf{발주사 (PMO)} & 형상관리 계획 CM-01(1.6.2 발주) & 일 1회 자동 동기화 설정, 클린 빌드 재현 시험 월 1회로 & P1 PMO 형상 담당 & 10-16 & 없음 & 완료 10-02 \\ \hline
AF-1-08 & 컨틴전시 집행 기록 중 R-03 0.50(8월)의 CCB 보고 누락 — 집행 절차(RACI 19행 CCB I 월 보고) 미준수 & 보통 & 발주사 & \textbf{발주사 (P1 PM)} & 집행 결재는 완료(P1 PM 08-12, 재경 확인 08-14), 8월 CCB(08-20) 안건 누락은 PMO 실수 & 10-15 CCB 소급 보고, 컨틴전시 집행 → CCB 안건 자동 등재 절차(집행 결재 시 안건 생성), 윤태호 사본 & P1 PM (PMO) & 10-15 & 없음 & 진행 (CCB 10-15) \\ \hline
\fi
\end{longtable}}
\B{\small\g{귀속 집계: 감리 기재 수행사 [ ]·발주사 [ ] → 발주사 판정 수행사 [ ]·발주사 [ ]·공동 [ ]. 등급: 중대 [ ]·보통 [ ]·경미 [ ]}\par}{\small 귀속 집계: 감리 기재 수행사 6·발주사 2 → \textbf{발주사 판정 수행사 3(02·05·06)·발주사 4(03·04·07·08)·공동 1(01)}. 등급: 중대 1·보통 4·경미 3.\par}
\landscapeoff

% ------------------------------------------------------------------ 6-1 이의 제기 서면 (세로)
\secbar[70mm]{6-1. 감리 지적 귀속 이의 제기 서면}
\guide{수신(감리원)·참조(스폰서·프로그램 PMO장·수행사 PM)·발신·일자·제목 / ① 지적 사실의 인정 범위 ② 귀속 재분류 요청과 근거(WBS 담당 열·RACI·기록 일자) ③ 발주사 조치 계획(담당·기한) ④ 요청 사항(정정·회신 기한·기준 합의) / 첨부 목록. 사실이 아니라 \textbf{귀속}에 대해서만 이의한다 — 사실을 다투면 감리와의 관계가 깨지고, 귀속을 다투지 않으면 검수·변경협의체 기록이 틀린다.}
\begin{fieldbox}\small
\ifwbblank
\g{수신} \hrulefill\quad \g{참조} \hrulefill\par
\g{발신} \hrulefill\quad \g{일자} \hrulefill\par
\g{제목} \hrulefill\par
\g{1. 지적 사실의 인정 — 대상 지적 번호, 인정 범위(사실은 인정하되 귀속만 다툰다는 것을 첫 문단에)}\lines{3}{7mm}
\g{2. 귀속 재분류 요청 — 지적별로: WBS 담당 열(발주/수행/공동) · RACI 행의 A · 결정 요청과 회신(일자) · 수행사의 의존 통지 기록(일자) · 수행사가 단독으로 확정할 권한이 없는 이유}\lines{6}{7mm}
\g{3. 발주사 조치 계획 — 담당(실명·직위) · 기한 · 시정 확인 회의 일자}\lines{3}{7mm}
\g{4. 요청 사항 — ① 정본 귀속란 정정 ② 회신 기한(검수 근거 사용 시점) ③ 향후 귀속 판정 기준 합의}\lines{3}{7mm}
\g{첨부: 1. \hrulefill\ 2. \hrulefill\ 3. \hrulefill}\par
\else
\textbf{수신} 정하늘 수석감리원 (감리법인) \quad \textbf{참조} 정미란 프로그램 스폰서, 프로그램 PMO장, 이재현 한빛디지털 PM\par
\textbf{발신} ONE ONDAM P1 PM \quad \textbf{일자} 2026-09-29\par
\textbf{제목} 중간감리 1차 지적사항 AF-1-03·AF-1-04의 귀속 재분류 요청\par
\textbf{1. 지적 사실의 인정.} 2026-09-25 접수한 중간감리 1차 수행 결과 보고서의 지적사항 8건 중 AF-1-03(P2 설비 인터페이스 규격 설비측 확인 서명 부재)과 AF-1-04(가맹 SW 배포 범위 설계 미확정)의 \textbf{지적 사실은 인정}하며, 두 항목이 G2·G3 통과에 필요한 산출물의 미완 상태임에 이견이 없습니다.\par
\textbf{2. 귀속 재분류 요청.} 두 지적은 보고서에 “수행사 설계 미흡”으로 기재되어 있으나, 아래 근거에 따라 원인은 \textbf{발주사(프로그램 계층 및 영업본부)의 결정 미완}이며, 귀속을 “발주사”로 재분류해 주실 것을 요청합니다.\par
\begin{itemize}
\item AF-1-03: 해당 규격(IF-P2-01~04)은 WBS 1.3.1 \textbf{발주사 담당} 작업패키지의 산출물이며(범위 기술서 v1.0 §8, 담당 “발주 박세연·한동석 협의”), P2 전달 책임은 RACI 9행에 따라 \textbf{프로그램 PMO장(A)}입니다. 규격은 T0 이전인 2026-04-17에 전달되었고(발송 기록 첨부 1), 설비 벤더의 05-29 변경 요청 2건에 대한 프로그램 PMO의 회신이 현재까지 미결입니다(회신 요청 06-05·07-10·08-21, 첨부 2). 수행사는 06-08·08-24 두 차례 미결 상태를 서면 통지했습니다(첨부 3). 수행사가 확인 서명 없이 규격을 확정할 권한은 없습니다.
\item AF-1-04: 비호환 단말 32점(8.2\%)의 대응안 결정은 WBS 1.1.6 \textbf{발주사 담당}이며 RACI 11행 \textbf{오정근 영업본부장(A)} 사항입니다. 대응안 3개는 04-30 조사 보고에 제시되었고(첨부 4), 결정은 가맹 부속서 협상 이후로 유보한다는 영업본부 회신(05-08, 첨부 5)이 있습니다. 배포 패키지 설계는 WBS 사전 1.1.6 카드의 외부 의존 조항에 따라 결정 이후 착수하도록 계획되어 있었습니다.
\end{itemize}
\textbf{3. 발주사 조치 계획.} AF-1-03은 프로그램 PMO장이 P2 PM과 협의해 2026-10-16까지 규격 확정 회신과 확인 서명을 완료하고, 수행사는 회신 후 5영업일 내 규격서 v1.1을 갱신합니다. AF-1-04는 오정근 영업본부장이 G2(2026-11-27)까지 대응안을 결정하며, 미동의 매장의 컷오버 후 2차 배포 옵션을 포함합니다. 두 건의 시정 확인은 11-06 감리 시정 확인 회의에서 받겠습니다.\par
\textbf{4. 요청 사항.} ① 감리 보고서 정본의 귀속란을 “발주사”로 정정, ② 정정본을 R2 검수 근거 문서로 사용할 수 있도록 10-08까지 회신, ③ 향후 감리 지적의 귀속 판정 시 WBS 담당 열(발주/수행/공동)과 RACI를 기준으로 하는 데 합의. 이 요청은 지적의 내용이나 등급에 대한 이의가 아니며, 검수·변경협의체에서 원인 귀속이 사실과 다르게 기록되는 것을 방지하기 위한 것입니다.\par
\textbf{첨부}: 1. 규격 전달 기록(04-17) 2. 회신 요청 기록 3건 3. 수행사 미결 통지 2건 4. 단말 조사 보고(04-30) 5. 영업본부 회신(05-08) 6. WBS 트리 담당 열·RACI 발췌\par
\fi
\end{fieldbox}

% ------------------------------------------------------------------ 7
\secbar[70mm]{7. RTM 커버리지 감리 대조}
\guide{감리 시점의 RTM 집계(층별 건수, 상위 연결·하위 분해·시험 연결·검수 연결률)와 G2 기준(설계·시험 연결 100\%) 대비 미달분의 성격 — “아직 설계 안 된 R3·R4 항목”인지 “누락”인지 — 과 도달 계획. Day2 ⑦ 항목 10 표에 시점 열을 붙인다. 감리가 미달을 지적하면 급히 빈 SyRS로 채우지 않는다.}
{\footnotesize\setlength{\tabcolsep}{2.5pt}
\begin{tabularx}{\linewidth}{|L{24mm}|C{9mm}|C{10mm}|C{11mm}|C{14mm}|C{12mm}|C{16mm}|Y|L{30mm}|}\hline
\hd{시점} & \hd{층} & \hd{건수} & \hd{상위 연결} & \hd{하위 분해} & \hd{시험 연결} & \hd{검수 연결} & \hd{미달 성격} & \hd{도달 계획} \\ \hline
\ifwbblank
\rd{9mm}\g{[전 게이트]} & StRS & & & & & & \g{[설계 미착수 / 릴리스 항목 / 누락]} & \\ \hline
\rd{9mm}\g{[G1]} & StRS & & & & & & & \\ \hline
\rd{9mm}\g{[감리 시점]} & StRS & & & & & & \g{[미분해 n = 어느 릴리스·HR, 누락 n]} & \g{[스프린트 종료 → 동결 → 게이트]} \\ \hline
\rd{9mm}〃 & SyRS & & & — & & — & & \\ \hline
\rd{9mm}〃 & SRS & & & — & & — & & \\ \hline
\rd{9mm}\g{[다음 게이트 목표]} & StRS & & & & & & & \\ \hline
\else
S6 종료 (Day2, 03-27) & StRS & 214 & 100\% & 71\% & 43\% & 0\% & 설계 미착수 & G1까지 100\% → 실제 G1 89\% \\ \hline
G1 (04-30) & StRS & 214 & 100\% & 89\% & 58\% & 0\% & R3·R4 항목 & — \\ \hline
\textbf{감리 1차 (09-18)} & StRS & 214 & 100\% & \textbf{88\%} (미분해 26) & \textbf{76\%} & 12\% (R1·R2 인수분) & 미분해 26 = R3 옴니 11·B2B EDI 6·R4 잔여 9 — 스프린트 내 설계 대상, \textbf{누락 0}. CR-06 승인 16건 중 4건은 FZ-06 반영으로 분모 증가 & S22 종료(11-08) 설계 확정 → FZ-11(11-25) → \textbf{G2 100\%} \\ \hline
〃 & SyRS & 399 & 100\% & — & 71\% & — & & \\ \hline
〃 & SRS & 575 & 100\% & — & 66\% & — & & \\ \hline
G2 목표 (11-27) & StRS & 215 (CR-08) & 100\% & 100\% & 100\% & 45\% (R3 인수분) & & G3 검수 연결 100\% \\ \hline
\fi
\end{tabularx}}

% ------------------------------------------------------------------ 8
\B{\clearpage}{}
\secbar[60mm]{8. 품질 허용범위 대비 상태와 에스컬레이션}
\guide{격자 품질 축(작업패키지·프로젝트)과 지속가능성 축(보안 High 미조치 0·개인정보 사고 0)의 현재 값·판정·에스컬레이션 기록(누가 언제 보고했고 P1 PM이 3영업일 내 무엇을 결정했는가), 감리 시정조치 이행 원칙(허용범위 무관), 다음 판정 시점. 품질 축을 G3에서 처음 판정하지 않는다.}
{\footnotesize\setlength{\tabcolsep}{2.5pt}
\begin{tabularx}{\linewidth}{|L{24mm}|L{30mm}|L{26mm}|L{16mm}|Y|L{22mm}|}\hline
\hd{축} & \hd{값 (데이터 일자)} & \hd{허용범위 (격자)} & \hd{판정} & \hd{에스컬레이션 (누가 · 언제 · 결정)} & \hd{다음 판정} \\ \hline
\ifwbblank
\rd{16mm}품질 — 작업패키지 & & & \g{[이내/경고/초과]} & \g{[보고자 → P1 PM 일자(인지+1영업일) / 결정 일자(3영업일) · 결정 내용]} & \\ \hline
\rd{12mm}품질 — 프로젝트 & \g{[미판정 시 전망값]} & & & & \\ \hline
\rd{12mm}지속가능성 & \g{[High n건 / 사고 n]} & & & & \\ \hline
\rd{10mm}감리 시정조치 & \g{[n건 중 완료·진행]} & 허용범위 무관 이행 & & & \\ \hline
\else
품질 — 작업패키지 & 0.132건/FP & 0.13 (수행사 PM) & \textbf{초과} & 1단계: 이재현 → P1 PM 보고 09-28(인지 09-25 + 1영업일). P1 PM 결정 10-01(3영업일): ① 결함 해소 전담 2명 10월 배치(R3 설계 병행 해제) ② backlog 10-31 60 이하 ③ 시험 노력(TC 실행) 주간 보고 추가 & 10-31 \\ \hline
품질 — 프로젝트 & 미판정 (통합시험 전) & 0.4 + 심각도 1·2 미해결 0(G3) & 전망 약 0.09 & — & 2026-12-07 착수 시 내부 기준(심각도 2 미해결 0) \\ \hline
지속가능성 & 보안 High 2건 조치 중(D-0966·VA-1-02), 개인정보 사고 0, 안전·환경 0 & High 미조치 0(G3), 사고 0 & \textbf{경고} (기한 내 조치 시 충족) & 병행: 최영주 보고 09-10, 재진단 10-23 & 10-23 재진단 \\ \hline
감리 시정조치 & 8건 중 완료 1(07)·기준서 완료 1(05)·진행 6 & 허용범위 무관 이행 & 이행 중 & 시정 확인 회의 11-06 & G2 감리 2차 \\ \hline
\fi
\end{tabularx}}

% ------------------------------------------------------------------ 9
\secbar[55mm]{9. 다음 조치와 통합시험 준비 연계}
\guide{결함 backlog 해소 계획(수치·기한), 심각도 2 기한, 감리 시정 확인 회의 일자·참석, 통합시험(A14) 착수 조건(심각도 2 미해결 0 · TC 확정 · 회귀 자동화 · RTM 시험 연결 100\%)과 판정일, 다음 품질 보고 시점. 날짜와 담당이 있는 5~7줄 — 품질 보고가 통합시험 계획과 따로 놀지 않게.}
\fieldunbreakable
\begin{fieldbox}\small
\ifwbblank
\g{backlog: [현재] → [기한]까지 [목표] ([인력·주당 해소])}\lines{1}{7mm}
\g{심각도 2: [n건] → [기한] 0 (선행 조건이 있는 건 명시)}\lines{1}{7mm}
\g{감리 시정 확인 회의: [일자] ([참석])}\lines{1}{7mm}
\g{통합시험 착수 조건과 판정일: }\lines{2}{7mm}
\g{측정 개선(열 추가 등): }\lines{1}{7mm}
\g{다음 품질 보고: [데이터 일자] ([부속 관계])}\lines{1}{7mm}
\else
\begin{itemize}
\item backlog 116 → 10-31까지 60 이하(전담 2명, 주 25건 해소), 심각도 2 7건 → 10-23 전 0 (D-0958은 정본 판정 착수 10-19 전 필수).
\item 감리 시정 확인 회의 11-06(정하늘·P1 PM·이재현·프로그램 PMO장·오정근 대리). AF-1-03 회신 10-16, AF-1-04 결정 11-27은 G2 감리 2차에서 확인.
\item 통합시험(A14) 착수 조건 12-04 판정: 심각도 2 미해결 0 / TC 400건 확정(Must 경로 우선 순서, R-19) / 회귀 자동화 도구 가동(EX-01 대안 1) / RTM 시험 연결 100\%.
\item 시험 노력(주당 실행 TC)을 결함표 열로 추가(10월부터) — 8월 착시 재발 방지.
\item 다음 품질 보고: 10-31 데이터(성과보고 제10호 부속).
\end{itemize}
\fi
\end{fieldbox}
\B{\small\g{서명: PMO 품질 담당 \hrulefill\ 수행사 PM \hrulefill\ P1 PM \hrulefill}\par}{\small 서명(10-05): P1 PMO 품질 담당 / 이재현(결함 데이터 대조) / P1 PM.\par}
\end{document}
